\section{4. Scaling: The Three-Axis Mental Model}
% ============================================================

\begin{bluf}
Modern AI progress isn't just ``bigger models.'' It's a Pareto optimization across three independent scaling axes. Anthropic's RL scaling team works on axis 2. Understanding why all three matter shows systems thinking.
\end{bluf}

\subsection{The Three Axes}

\textbf{Axis 1: Pretraining compute.} More parameters $\times$ more data = better base model. Power-law improvements: each 10$\times$ compute gives diminishing but substantial gains. Chinchilla's key insight: scale parameters and data \textit{equally}. The original Kaplan work over-allocated to parameters.

\textbf{Axis 2: RL compute.} This is the RL scaling team's domain. RL training on top of the pretrained model. DeepSeek R1 showed that even \textit{modest} RL compute (much less than pretraining) can produce dramatic capability jumps. The constraint isn't compute --- it's having clean enough feedback signals to spend that compute on.

\textbf{Axis 3: Test-time compute.} More thinking time at inference: chain-of-thought, search, self-verification, retries. This is what makes reasoning models ``slow but smart.''

\begin{sayit}[Scaling]
``The interesting thing about the current moment is that we're not just scaling one dimension anymore. Pretraining, RL training, and test-time compute are three separate levers, and the question is how to allocate across them. Recent work has shown that RL compute is dramatically under-exploited relative to pretraining --- modest RL budgets produce huge capability jumps. The bottleneck isn't the compute; it's having verification signals clean enough to spend that compute on.''
\end{sayit}

\subsection{The Compute vs.\ Scaffolding Tradeoff}

The key framing: ``How much am I willing to burn compute, versus how much am I willing to burn dollars on people's time to give scaffolding or give rewards?''

\textbf{Key tension:} You can invest in better reward engineering (human effort, clever verification) or just throw more RL compute at imperfect rewards. But ``things which can help at smaller scales can actually hurt at larger scales'' --- a reward signal that's clean enough for 1B parameters might be too noisy for 70B.

\begin{deeper}
This is why RL scaling teams validate algorithmic ideas at small scale before betting big: ``You want to be sure that you've algorithmically got the right thing, and then when you bet and you do the large compute spend on the run, then it'll actually pay off.'' The risk of a large RL training run is that you don't know if your reward signal holds up at scale until you've spent the compute.
\end{deeper}

% ============================================================
\section{5. Anthropic's Approach: Constitutional AI}
% ============================================================

\begin{bluf}
Anthropic's signature innovation is using AI feedback instead of human feedback (RLAIF). You should understand it because it's the company's philosophical foundation and it connects to scalable oversight --- the question of how to supervise systems smarter than you.
\end{bluf}

\textbf{The core idea:} Instead of paying humans to rank outputs, you write a set of principles (a ``constitution'') and have the AI critique and revise its own outputs against those principles. Then you train a preference model from the AI's judgments, not human ones.

\textbf{Why it matters:}
\begin{itemize}
  \item Human feedback doesn't scale and is expensive
  \item For safety-critical behaviors, you need a feedback loop that works 24/7
  \item If ``evaluating is easier than generating,'' then AI can provide reliable feedback even on tasks where it can't yet perform well
\end{itemize}

\textbf{The scalable oversight question:} How do you supervise a model that's smarter than your evaluators? Constitutional AI is Anthropic's bet: formalize the criteria, automate the feedback loop, and scale oversight through AI-assisted evaluation.

\begin{sayit}[Constitutional AI]
``What I find interesting about Constitutional AI is that it's really a scalable oversight play. Human feedback is expensive and doesn't scale. If you can get AI to reliably evaluate against explicit principles, you've broken through the human bottleneck. The question is whether the AI evaluator is robust enough --- and that loops back to verification quality.''
\end{sayit}

% ============================================================
\section{5.5. The Verification Bottleneck Pattern}
% ============================================================

\begin{bluf}
Every frontier problem in RLVR reduces to a verification problem. If you can verify, RLVR works. If you can't, you're stuck with RLHF's gameable proxies. The skill that matters most is designing verifiers.
\end{bluf}

\subsection{The Pattern}

``Why is robotic manipulation harder than locomotion? Because it's hard to verify correctness.'' This pattern appears everywhere:

\begin{itemize}
  \item \textbf{RL training:} Reward hacking = verification failure --- the model scored well on something your verifier didn't catch
  \item \textbf{Performance engineering:} How do you verify a kernel is correct, not just fast?
  \item \textbf{Hardware evaluation:} How do you verify benchmarks extrapolate to real workloads?
  \item \textbf{Scaling:} How do you verify small-scale results predict large-scale outcomes?
\end{itemize}

\begin{sayit}[The verification pattern]
``The thing I keep coming back to is that every hard problem in RLVR is secretly a verification problem. Manipulation is harder than locomotion: you can tell if a robot walked, but did it actually pick up the object correctly? The same pattern runs through all of it --- reward hacking is a verification failure, slow PE is a verification failure (fast but wrong kernels), and bad hardware benchmarks are a verification failure. Once you see it as one problem, you start asking: what's the hardest thing to verify in this system, and how do we build a verifier for it?''
\end{sayit}

% ============================================================
